% Study 4: supplementary methods, the sample and its relation to the CogSci 2025 count, the relation
% to the CogSci 2025 analysis, and the results that no other part of the paper prints. Numbers come
% only from macros in generated/numbers.tex.
\label{sec:suppB-study4}

Study~4 asked whether two things said about an original, its label and the hours Alex spent on it, carry over to ratings of the second works he was said to have made from it with his design and sketching tools. Its data were collected from 27 January to 2 February 2025, and our CogSci 2025 paper\cite{kamper2025originalsin} reported it as its Study~2. Its registration, a dated record of its hypotheses and analysis plan, was filed on the Open Science Framework (OSF) on 8 October 2026, after analyses of the same data had been reported in that paper and in a September 2025 talk (Supplementary Table~\ref{tab:supp-verdicts}). The main text and Supplementary Table~\ref{tab:supp-labeleffects} report the label's effect, and Supplementary Note~6 discusses its registered tests at each level. This section gives the materials, procedure and sample, explains why the registered sample is smaller than the one the CogSci 2025 paper gives, sets out how the two analyses differ, and reports the results that no other part of the paper prints.

\subsection*{Materials}

Study~4 used the 30 images of Study~2: ten originals, all generated with OpenAI's DALL·E, an AI image generator, each with a Slight and a Dramatic version that we made by hand and a None second work, the original shown again unchanged (Section~\ref{sec:suppA-registrations}; Supplementary Note~4). The Human label was therefore a cover story. Each participant saw six of the ten originals, three with a human or humanoid figure and three without, each with a stated time and each followed by its three second works.

The Human label read: `There is a digital artist named Alex. Alex is a skilled artist who creates digital artworks. Alex uses his graphic design skills and digital sketching tools.' The AI label, the same in all five studies, read: `There is a digital artist named Alex. Alex is a skilled artist who creates digital artworks. Alex develops and inputs prompts into an artificial intelligence (AI) system. The AI system generates digital artworks based on those prompts.'

Before each original, the survey inserted a stated time, `less than 1 hour', `10 hours' or `100 hours', where \textit{time} appears below. The Human group read `Alex had a new idea for a digital artwork. Alex spent \textit{time} working with his graphic design and digital sketching tools, revising as his work as he goes, which creates a digital artwork.' The AI group read `Alex had a new idea for a work of digital artwork. Alex spent \textit{time} developing and inputting prompts into the AI image generator, revising each prompt as he goes, which creates a digital artwork.' Both went on: `Click the arrow to see the digital artwork created. Note that you will be required to stay on the page for at least 8 seconds before proceeding.'

Both groups read the same description of each second work. Before the first of the three, it was `Alex looked at the first artwork. He decided to produced a secondary digital artwork. Using the first artwork as a base, he produced a second digital artwork with his own graphic design skills and digital sketching tools.' Before each of the other two, it was `Later, Alex looked at the very first artwork again. Alex decided to produce a different secondary digital artwork. Using the very first artwork as a base, he produced a different secondary digital artwork with his own graphic design skills and digital sketching tools.' Each ended `Click the arrow to see the digital artworks:'. The quotations reproduce the survey exactly, slips included. No time was stated for the second work.

The three ratings were those of Study~2, with two changes: the essence and creativity items now carried definitions. The essence item read `How much has the essence of the first work changed in the second work? Essence refers to the core sense and feeling that one has while observing a work of art that makes that work unique.' The creativity item read `How much creativity was involved in the production of the second work? Creativity is defined as the presence of at least a minimal original, intellectual contribution to the creation of the work.' The transformativeness item and every answer scale were as in Study~2, scored 1 to 5 by answer position. The rating screen was also laid out as in Study~2, with the three ratings first and a required explanation box for each below them, and a script again blocked copying and pasting in 21 of the survey's 90 explanation boxes and not in the others. Two transformativeness items, both for one original, again had six answer points; in Study~4 they kept the usual question and definition. As registered, their \SFourNSixPointDropped{} ratings were left out of every model. The slips in three essence answer labels and one creativity question recurred; scoring by answer position is unaffected.

Study~4 thus differs from Study~2 in its design and in its wording. In design, it showed six originals per person rather than ten, each with a stated time. In wording, its essence and creativity items carried definitions, its Human label and vignettes followed Study~1's, its AI vignette mentioned revised prompts, and its second works were said to be made with Alex's `own graphic design skills'.

\subsection*{Procedure}

Participants took part online in a Qualtrics survey. By our account, participants were students in the UCLA Psychology subject pool who took part for course credit. Every analyzed record carries the code that the survey sent, at its end, to the pool's page for granting credit. The consent form, the same in all five studies, also offered hourly pay to people who signed up outside the pool, and no record shows that anyone took that route. The survey assigned each participant at random to one label for the whole session, presenting each label equally often. It also set a quota of 62 participants for each label group, a planned number that neither stopped anyone nor moved anyone to the other label. A script attached a stated time to each original, drawing at random among the times not yet used twice, so each participant saw two originals at each time in a random order. Each original was shown alone for at least 8 seconds, and each second work beside it for at least 10 seconds. The order of the three levels was random and was not recorded. Each participant rated 18 second works, giving 54 ratings and 54 explanations.

The task question, the arithmetic question (after the trials) and the closing screen, which asked the seriousness question before the background questions, were as in Study~2. No question asked participants to recall the label (a manipulation check). The debriefing was the same as in the other four studies; Section~\ref{sec:suppA-study1} quotes it and corrects its statement of the law. Participants read it only after their last rating.

\subsection*{Exclusions and sample}

The registered exclusion rules, applied in order, left \SFourN{} of the \SFourNRecords{} records in the survey export. First, \SFourNPreviewRemoved{} preview records, which are test runs, were removed. Then \SFourNUnfinishedRemoved{} unfinished records were removed, and \SFourNNotSeriousRemoved{} participants who gave any answer but `Yes, I did' to the seriousness question, a blank included. No one answered the arithmetic question wrongly. Last, the registration defines a valid completion as one that rates exactly six originals, each named once in the survey's six image slots, two at each stated time. Of the remaining records, \SFourNInvalidRemoved{} failed this rule: one had rated all ten originals, and one named the same original in two slots, so its stated times were not split two, two and two. Of the \SFourN{} participants analyzed, \SFourNHuman{} saw the Human label and \SFourNAI{} the AI label. They gave \SFourNRatings{} analyzed ratings and \SFourNJustifications{} explanations.

The CogSci 2025 paper gives this study's sample as \SFourNBeforeValidity{} participants. Its count before the seriousness question, and the number it removed at that question, both match the pipeline's, so its reported sample is the count left before the validity rule, and the difference from the registered sample is the \SFourNInvalidRemoved{} records that the rule removes. The paper starts from a smaller count of recruited participants than the export's records that are not test runs, and it removes correspondingly fewer unfinished sessions. Study~4's registration adds that the analyses of variance behind the CogSci 2025 paper were fitted to as many participants as the registered sample holds, all but one of them in that sample. The degrees of freedom that the CogSci 2025 paper gives for this study fit that number, not the sample it states.

The participants' mean age was \SFourAgeMean{} years (standard deviation \SFourAgeSD; range \SFourAgeMin{} to \SFourAgeMax). Of the \SFourN, \SFourNWomen{} described their gender as female, \SFourNMen{} as male, and \SFourNGenderOther{} chose `Other' or `Prefer not to say' or gave no answer. The answers on race, native language and formal artistic training are not summarized here. For this study, the CogSci 2025 paper gives age and gender figures for the participants it recruited, not for the analyzed sample.

\subsection*{How the registered analysis differs from the CogSci 2025 analysis}

The CogSci 2025 paper analyzed these data with mixed analyses of variance (ANOVAs). A mixed ANOVA compares condition means computed from each participant's average rating in each condition; here the label varied between participants, and the stated time and the level within them. The registered analysis instead fits a linear mixed model to every rating, a regression that gives each participant and each original its own baseline. It fixes fifteen confirmatory tests, five per measure: two steps of modification (H1), an equivalence test of the label (H2) and two $F$~tests of the stated time (H3). It corrects the fifteen together with Holm's method, which raises each $p$-value so that the chance of calling any of them significant by mistake stays at .05. A $p$-value is the chance of a result at least this large if there were no effect. An equivalence test asks whether an effect is shown to lie inside a stated bound, here $\pm$\SourceBound{} points, half of Study~1's creativity penalty on originals. An $F$~test asks whether a factor changes the mean rating.

The two analyses agree that ratings rose steeply with the level and that the label's overall effect was not significant; in the registered model that $F$~test lies outside the registered set and is context only (Supplementary Table~\ref{tab:supp-labeleffects}). They part on the hours stated for the original. The CogSci 2025 paper found an effect of those hours on creativity, with the second works of 100-hour originals rated lowest. In the registered model the $F$~test of stated hours on creativity was not significant (uncorrected $p$~\SFourCreatSourceTimeP, Holm $p$~\SFourCreatSourceTimePHolm), the ordinal check agreed, and the predicted creativity mean at 100 hours minus the mean at under 1 hour was \SFourCreatTimeContrastEst{} points, inside the bound (registered secondary; Supplementary Table~\ref{tab:suppB-s4}c). The analyses also differ in what they hold constant. Each participant saw a random six of the ten originals, each at a randomly assigned time, so a participant's average at one stated time also reflects which originals happened to receive it. A baseline for each original, as in the registered model, removes that source of variation; an ANOVA on participants' averages does not. Whether this difference accounts for the different results was not tested.

The CogSci 2025 paper also found that the label's effect on essence change differed by level, with a difference at Slight only. The registered secondary tests at each level likewise found differences favoring the Human label at Slight, on transformativeness and essence change, without correction (Supplementary Table~\ref{tab:supp-labeleffects}). Outside the registered set of tests, the model also has a term that lets the label's effect differ by level. Examined as an exploratory, uncorrected test, it was significant on transformativeness ($F$(\SFourTransLabelByModDfNum,~\SFourTransLabelByModDfDen)~=~\SFourTransLabelByModF, $p$~\SFourTransLabelByModP) and essence change ($F$(\SFourEssLabelByModDfNum,~\SFourEssLabelByModDfDen)~=~\SFourEssLabelByModF, $p$~\SFourEssLabelByModP), and not on creativity ($p$~\SFourCreatLabelByModP).

Study~4 was not registered with any registry before data collection. Its only planning document was a form that was never submitted and was last edited during collection. That form planned analyses of variance and a measure of perceived effort that the survey did not contain, and it predicted that stated time would raise ratings most under the AI label. These data do not show that pattern: the registered secondary test of whether the stated hours changed the label's effect was not significant on any measure (Supplementary Table~\ref{tab:suppB-s4}c).

\subsection*{Results not reported elsewhere}

Each result below carries its status. A registered confirmatory test is one that the registration names as confirmatory, judged by its rule after Holm's correction; a result is called significant when its $p$-value is below .05. A registered secondary analysis is one the registration names outside that set, judged without correction. An exploratory analysis was named as exploratory or not named at all, and carries no verdict.

\paragraph*{Registered tests and predicted means.}
Supplementary Table~\ref{tab:suppB-s4} gives every registered test of H1 and H3 with its ordinal check, the registered secondary tests of the stated hours, and the predicted means. The ordinal check is a cumulative-link mixed model, which treats the five answers as ordered categories rather than evenly spaced numbers; it reached the same verdict as the linear model on all twelve tests it repeated. All three registered linear models and all three ordinal models converged, meaning that the fitting procedure settled on a solution, and none was singular, that is, none estimated a source of variation at zero.

% Study 4: the registered tests of H1 and H3 with the ordinal check, the registered secondary
% tests of the hours stated for the original, and the predicted means.
% Every number is a macro from generated/numbers.tex; supplement/tables/README_suppB.md lists
% the pipeline tables behind them.
\begin{table}[htbp]
\centering
\footnotesize
\caption{\textbf{Study~4: the registered tests of modification (H1) and of the hours stated for the original (H3), and the predicted means.} Ratings are on 1-to-5 scales. \textbf{a},~H1, registered confirmatory: the difference between predicted means at adjacent levels, averaged over label and stated hours. \textbf{b},~H3, registered confirmatory: $F$~tests of whether the hours stated for the original, alone or by level, change the mean rating. Panels a and b are corrected together with H2 by Holm's method across fifteen tests. \textbf{c},~Registered secondary tests, each uncorrected: whether the stated hours change the label's effect, and 100 hours minus under 1 hour, read with an equivalence test against $\pm$\SourceBound{} points. \textbf{d},~Predicted means (descriptive).}\label{tab:suppB-s4}
\setlength{\tabcolsep}{4pt}
\begin{tabular}{@{}lcccccc@{}}
\multicolumn{7}{@{}l}{\textbf{a}\enspace H1: ratings rise with the level}\\
\toprule
 & & & & & \multicolumn{2}{c}{Ordinal check}\\
\cmidrule(l){6-7}
Step & Estimate & SE & $t$ ($df$) & After Holm & Log-odds & After Holm\\
\midrule
\multicolumn{7}{@{}l}{\textit{Transformativeness}}\\
Slight minus None & \SFourTransSlightMinusNoneEst & \SFourTransSlightMinusNoneSE & \SFourTransSlightMinusNoneT{} (\SFourTransSlightMinusNoneDf) & $p$~\SFourTransSlightMinusNonePHolm & \SFourTransSlightMinusNoneOrdEst & $p$~\SFourTransSlightMinusNoneOrdPHolm\\
Dramatic minus Slight & \SFourTransDramaticMinusSlightEst & \SFourTransDramaticMinusSlightSE & \SFourTransDramaticMinusSlightT{} (\SFourTransDramaticMinusSlightDf) & $p$~\SFourTransDramaticMinusSlightPHolm & \SFourTransDramaticMinusSlightOrdEst & $p$~\SFourTransDramaticMinusSlightOrdPHolm\\
\addlinespace[3pt]
\multicolumn{7}{@{}l}{\textit{Essence change}}\\
Slight minus None & \SFourEssSlightMinusNoneEst & \SFourEssSlightMinusNoneSE & \SFourEssSlightMinusNoneT{} (\SFourEssSlightMinusNoneDf) & $p$~\SFourEssSlightMinusNonePHolm & \SFourEssSlightMinusNoneOrdEst & $p$~\SFourEssSlightMinusNoneOrdPHolm\\
Dramatic minus Slight & \SFourEssDramaticMinusSlightEst & \SFourEssDramaticMinusSlightSE & \SFourEssDramaticMinusSlightT{} (\SFourEssDramaticMinusSlightDf) & $p$~\SFourEssDramaticMinusSlightPHolm & \SFourEssDramaticMinusSlightOrdEst & $p$~\SFourEssDramaticMinusSlightOrdPHolm\\
\addlinespace[3pt]
\multicolumn{7}{@{}l}{\textit{Creativity}}\\
Slight minus None & \SFourCreatSlightMinusNoneEst & \SFourCreatSlightMinusNoneSE & \SFourCreatSlightMinusNoneT{} (\SFourCreatSlightMinusNoneDf) & $p$~\SFourCreatSlightMinusNonePHolm & \SFourCreatSlightMinusNoneOrdEst & $p$~\SFourCreatSlightMinusNoneOrdPHolm\\
Dramatic minus Slight & \SFourCreatDramaticMinusSlightEst & \SFourCreatDramaticMinusSlightSE & \SFourCreatDramaticMinusSlightT{} (\SFourCreatDramaticMinusSlightDf) & $p$~\SFourCreatDramaticMinusSlightPHolm & \SFourCreatDramaticMinusSlightOrdEst & $p$~\SFourCreatDramaticMinusSlightOrdPHolm\\
\bottomrule
\end{tabular}

\medskip
\begin{tabular}{@{}lccccc@{}}
\multicolumn{6}{@{}l}{\textbf{b}\enspace H3: the hours stated for the original}\\
\toprule
 & & & & \multicolumn{2}{c}{Ordinal check}\\
\cmidrule(l){5-6}
Term & $F$ ($df$) & Uncorrected & After Holm & $\chi^2$ ($df$) & After Holm\\
\midrule
\multicolumn{6}{@{}l}{\textit{Transformativeness}}\\
Stated hours & \SFourTransSourceTimeF{} (\SFourTransSourceTimeDfNum, \SFourTransSourceTimeDfDen) & $p$~\SFourTransSourceTimeP & $p$~\SFourTransSourceTimePHolm & \SFourTransSourceTimeOrdChisq{} (\SFourTransSourceTimeOrdDf) & $p$~\SFourTransSourceTimeOrdPHolm\\
Stated hours by level & \SFourTransSourceTimeByModF{} (\SFourTransSourceTimeByModDfNum, \SFourTransSourceTimeByModDfDen) & $p$~\SFourTransSourceTimeByModP & $p$~\SFourTransSourceTimeByModPHolm & \SFourTransSourceTimeByModOrdChisq{} (\SFourTransSourceTimeByModOrdDf) & $p$~\SFourTransSourceTimeByModOrdPHolm\\
\addlinespace[3pt]
\multicolumn{6}{@{}l}{\textit{Essence change}}\\
Stated hours & \SFourEssSourceTimeF{} (\SFourEssSourceTimeDfNum, \SFourEssSourceTimeDfDen) & $p$~\SFourEssSourceTimeP & $p$~\SFourEssSourceTimePHolm & \SFourEssSourceTimeOrdChisq{} (\SFourEssSourceTimeOrdDf) & $p$~\SFourEssSourceTimeOrdPHolm\\
Stated hours by level & \SFourEssSourceTimeByModF{} (\SFourEssSourceTimeByModDfNum, \SFourEssSourceTimeByModDfDen) & $p$~\SFourEssSourceTimeByModP & $p$~\SFourEssSourceTimeByModPHolm & \SFourEssSourceTimeByModOrdChisq{} (\SFourEssSourceTimeByModOrdDf) & $p$~\SFourEssSourceTimeByModOrdPHolm\\
\addlinespace[3pt]
\multicolumn{6}{@{}l}{\textit{Creativity}}\\
Stated hours & \SFourCreatSourceTimeF{} (\SFourCreatSourceTimeDfNum, \SFourCreatSourceTimeDfDen) & $p$~\SFourCreatSourceTimeP & $p$~\SFourCreatSourceTimePHolm & \SFourCreatSourceTimeOrdChisq{} (\SFourCreatSourceTimeOrdDf) & $p$~\SFourCreatSourceTimeOrdPHolm\\
Stated hours by level & \SFourCreatSourceTimeByModF{} (\SFourCreatSourceTimeByModDfNum, \SFourCreatSourceTimeByModDfDen) & $p$~\SFourCreatSourceTimeByModP & $p$~\SFourCreatSourceTimeByModPHolm & \SFourCreatSourceTimeByModOrdChisq{} (\SFourCreatSourceTimeByModOrdDf) & $p$~\SFourCreatSourceTimeByModOrdPHolm\\
\bottomrule
\end{tabular}

\medskip
\begin{tabular}{@{}lccccc@{}}
\multicolumn{6}{@{}l}{\textbf{c}\enspace Registered secondary tests of the hours stated for the original}\\
\toprule
 & \multicolumn{2}{c}{Label by stated hours} & \multicolumn{3}{c}{100 hours minus under 1 hour}\\
\cmidrule(lr){2-3}\cmidrule(l){4-6}
Measure & $F$ & Test & Estimate & 90\% CI & Equivalence test\\
\midrule
Transformativeness & \SFourTransLabelBySourceTimeF & $p$~\SFourTransLabelBySourceTimeP & \SFourTransTimeContrastEst & \SFourTransTimeContrastLoNinety{} to \SFourTransTimeContrastHiNinety & $p$~\SFourTransTimeContrastPTost\\
Essence change & \SFourEssLabelBySourceTimeF & $p$~\SFourEssLabelBySourceTimeP & \SFourEssTimeContrastEst & \SFourEssTimeContrastLoNinety{} to \SFourEssTimeContrastHiNinety & $p$~\SFourEssTimeContrastPTost\\
Creativity & \SFourCreatLabelBySourceTimeF & $p$~\SFourCreatLabelBySourceTimeP & \SFourCreatTimeContrastEst & \SFourCreatTimeContrastLoNinety{} to \SFourCreatTimeContrastHiNinety & $p$~\SFourCreatTimeContrastPTost\\
\bottomrule
\end{tabular}

\medskip
\begin{tabular}{@{}lccccc@{}}
\multicolumn{6}{@{}l}{\textbf{d}\enspace Predicted means}\\
\toprule
 & \multicolumn{3}{c}{By level} & \multicolumn{2}{c}{By label}\\
\cmidrule(lr){2-4}\cmidrule(l){5-6}
Measure & None & Slight & Dramatic & Human label & AI label\\
\midrule
Transformativeness & \SFourTransMeanNone & \SFourTransMeanSlight & \SFourTransMeanDramatic & \SFourTransMeanHuman & \SFourTransMeanAI\\
Essence change & \SFourEssMeanNone & \SFourEssMeanSlight & \SFourEssMeanDramatic & \SFourEssMeanHuman & \SFourEssMeanAI\\
Creativity & \SFourCreatMeanNone & \SFourCreatMeanSlight & \SFourCreatMeanDramatic & \SFourCreatMeanHuman & \SFourCreatMeanAI\\
\bottomrule
\end{tabular}

\smallskip
\begin{minipage}{\textwidth}
\footnotesize\raggedright
Estimates come from the registered linear mixed model, which gives each participant and each original its own baseline. SE: standard error. $df$: degrees of freedom by Satterthwaite's approximation. The ordinal check is a cumulative-link mixed model, which treats the five answers as ordered categories; it repeats the H1 and H3 tests as Wald tests and, having no rating-point scale, does not repeat the equivalence tests. Its H1 estimates are differences in the log-odds of a higher answer. The label-by-hours $F$~tests have two numerator degrees of freedom. Panel~d averages each level over label and stated hours, and each label over level and stated hours. Supplementary Table~\ref{tab:supp-labeleffects} gives the label's effect averaged over levels and at each level.
\end{minipage}
\end{table}


\paragraph*{Participants who wrote that a machine made the work.}
A fixed word list, run without access to labels or ratings, flagged explanations stating that a machine made the work. Of the \SFourMachineMadeNHuman{} Human-label participants, \SFourMachineMadeN{} wrote this at least once, in \SFourMachineMadeResponses{} of their \SFourMachineMadeNResponses{} explanations. This is a count of statements, not of disbelief. Without them, the registered secondary refit used \SFourNRefit{} participants and changed \SFourNRefitVerdictsChanged{} of the \SFourNRefitVerdicts{} registered verdicts, the essence-change equivalence test (main text). Every refit converged.

\paragraph*{Registered text checks.}
Two checks on the creativity explanations were registered for every study, and Study~4 ran both (registered secondary; Section~\ref{sec:suppA-text} describes both). In the first, ridge regression, a regression that shrinks each word's weight toward zero, predicted each rating and each participant's label from the words of explanations written by participants left out of the fit. The words predicted the rating: predicted and actual ratings correlated at \SFourTextRecRatingCreat{} (95\% confidence interval, CI, the range of values the data do not rule out at the 95\% level: \SFourTextRecRatingCreatLo{} to \SFourTextRecRatingCreatHi), and in runs with the ratings shuffled, 95\% of correlations stayed below \SFourTextRecRatingCreatNull. The words did not show the label. The AUC, the chance that a randomly chosen AI-label explanation scores above a randomly chosen Human-label one, so that 0.5 is chance, was \SFourTextRecArmCreat{} (\SFourTextRecArmCreatLo{} to \SFourTextRecArmCreatHi), an interval that reaches chance; the rating alone gave \SFourTextRecArmFromRatingCreat. Adding the words to the rating changed the AUC by \SFourTextRecIncrCreat{} (\SFourTextRecIncrCreatLo{} to \SFourTextRecIncrCreatHi), which is inconclusive. The second check predicted that AI-label participants would reuse their own words across explanations more than Human-label participants, measured as a person's word overlap among their own explanations minus their overlap with other people's explanations of the same second work, with explanation length held constant. The difference, AI minus Human, was \SFourTextOverlapCreatDiff{} (\SFourTextOverlapCreatLo{} to \SFourTextOverlapCreatHi; Holm $p$~\SFourTextOverlapCreatPHolm), so the prediction is inconclusive, which is not evidence that the groups reuse words equally. Every explanation was written after its rating, so these checks show what the words track, not why anyone rated as they did.

\paragraph*{Exploratory text analyses.}
The registration names the remaining text analyses as exploratory. No topic found without labels, a group of words that tend to occur together, carried more weight under one label than the other after a Benjamini--Hochberg correction, which limits the expected share of false positives among results called significant. Word tone and concreteness, scored from published norms, did not differ by label after the same correction. In predicting the rating, words such as `completely', `different' and `style' went with higher ratings and words such as `exact', `identical' and `copy' with lower ones, on all three measures. Their estimates are in the Study~4 tables that the analysis code writes (main text, Code availability). Section~\ref{sec:suppB-crossstudy} reports the label's effect at each stated time and level.

\clearpage
